% Getting Started, The Basics (PUB-REF-0001 AN-21 F-WPC: Registration,
% Installation, The Template, Starting InitechWord, Keys to Know, Function
% Keys, Pull-Down Menus, Menu Options, Codes, Exit; model: WordPerfect 5.1
% Reference pp. 3-11, docs/pubs/originals/wordperfect-51.md s3, s4.4).
% Law 1: os/iword/app/face_a.c (the screen, the status line, the keys, the
% File menu, the exit guard), spec/flair_iword_traces.mk (the opening of
% WORD.EXE from the Finder), the Makefile flair-data-rules (APPS\WORD.EXE on
% the shipped data volume), harness/pubs/iword_probe.c (which keys act).
% GATES: test-pubs-iword test-iword-face test-flair-iword test-pubs-flair-screens
\Gates{test-pubs-iword test-iword-face test-flair-iword test-pubs-flair-screens}
\WPHead{The Basics}
This section tells you what you need to know to begin working with
InitechWord: how the program is supplied, how to start it, the keys you use
most, and how the function keys, the menus and the hidden codes work.  If
you have used a word processor before, read it quickly and then try
\textit{A Brief Lesson} at the end of the section.  If you are new to word
processing, read it carefully and then work through the lessons of the
\textit{InitechWord Workbook}, which build on what is said here.

\WPSide{Registration}
Fill in the Warranty Registration Card that comes in the InitechWord
package and mail it to Initech Systems Corporation.  Registered users are
told of new editions of the program and of these publications, and may call
the Initech Help Desk for assistance (see \textit{Getting Help} below).
\draft{PUB-REF-0001 PK-13: the Warranty Registration Card is not yet
produced}

\WPSide{Installation}
InitechWord is supplied installed.  The program is the file
\Scr{WORD.EXE} in the \Scr{APPS} folder of the InitechOS data volume, the
disk that appears on the FLAIR desktop beside the system disk.  There is
nothing to copy and no installation program to run: when InitechOS is
running and FLAIR is on the screen, InitechWord is ready to start.

InitechWord keeps the documents you save in the top level of the same data
volume.  It shows that volume as drive \Scr{A:}, so a document you save
under the name \Scr{MEMO} appears on the status line as
\Scr{A:\wbs MEMO}.

The \textit{InitechWord Installation} card that comes in the package
repeats these points in brief and takes you through starting the program
the first time.

\WPSide{The Template}
The function-key template that comes with InitechWord lies above the
function keys of your keyboard.  It names the feature each function key
starts, alone and with the Shift, Alt and Ctrl keys.  On the template, each
key has up to four lines:
\begin{itemize}
\item the feature started by pressing Ctrl and the key, in red;
\item the feature started by pressing Alt and the key, in blue;
\item the feature started by pressing Shift and the key, in green;
\item the feature started by pressing the key alone, in black.
\end{itemize}
A line printed in grey names a WordPerfect feature InitechWord does not
carry out (see \textit{Function Keys} below), and a line is left blank where
the combination does nothing.  The quick list at the left of the template
names the other keys you use most often.

\WPSide{Starting InitechWord}
InitechWord is a FLAIR application.  To start it,
\begin{WPSteps}
\item Double-click the icon of the data volume on the FLAIR desktop.  The
volume opens in a window.
\item Double-click the \Scr{APPS} folder.  The folder opens in a window.
\item Double-click the \Scr{WORD.EXE} icon.
\end{WPSteps}
The InitechWord window opens over the desktop, below the menu bar, with
its title \Scr{InitechWord} in the title bar.  The menu bar shows a second
row of menus below the desktop's: the \textbf{File} menu of InitechWord
(see \textit{Pull-Down Menus} below).

\FlairShot{editing}{640}{440}{The InitechWord window: a short memo typed into a new document}

The window shows the editing screen.  Text appears in white on a blue
field, and nothing but your text is shown until you ask for something.  The
cursor is the cell drawn in reverse, white on blue: it marks where the next
character you type will go.

The bottom line of the window is the status line.  At its right it tells
you where the cursor is:
\begin{description}
\item[Doc] the document you are working on (always \Scr{Doc 1})
\item[Pg] the page the cursor is on
\item[Ln] how far the line the cursor is on lies from the top edge of the
page, in inches
\item[Pos] how far the cursor lies from the left edge of the paper, in inches
\end{description}
On a new document the cursor is at \Scr{Ln 1"} and \Scr{Pos 1"}: one inch
from the top edge and one inch from the left edge, inside the standard one
inch margins.  Once you have saved a document or retrieved one, its name is
shown at the left of the status line.  The same line is used to ask you
questions and to show menus, as described below.  See \InRef{Status Line}.

\WPSide{Keys to Know}
You type into InitechWord as you would on a typewriter, with these
differences.

\Topic{Enter}
Press \K{Enter} only to end a paragraph, to end a short line such as a
heading or a line of an address, or to make a blank line.  When a line
fills, InitechWord moves the word that does not fit to the next line for
you (see \InRef{Word Wrap}).  Each time you press \K{Enter}, InitechWord
puts a hard return code in the document.

\Topic{Tab}
\K{Tab} moves the text that follows to the next tab stop.  The stops are
every half inch.  See \InRef{Tab}.

\Topic{Backspace and Delete}
\K{Backspace} erases the character to the left of the cursor; \K{Delete}
erases the character at the cursor.  See \InRef{Delete}.

\Topic{Cursor Keys}
The arrow keys move the cursor through the text one character or one line
at a time.  \K{End} moves it to the end of the line.  \K{Home} is used in
front of the arrow keys to move farther: \K{Home, Left} goes to the start
of the line, and \K{Home, Home, Up} to the top of the document.  See
\InRef{Cursor Movement}.

\Topic{Cancel}
\FK{Cancel}{F1} backs out of any menu, prompt or question and returns you
to the document.  \K{Esc} does the same.  On the editing screen, with
nothing to back out of, Cancel offers to restore the text you deleted last.
See \InRef{Cancel} and \InRef{Undelete}.

\Topic{Shift, Alt and Ctrl}
These keys change what a function key does, as the template shows; with
\K{Enter}, Ctrl gives a hard page (\K{Ctrl-Enter}).  With every other key
InitechWord acts as if they were not held: Shift-Backspace erases as
Backspace does, Ctrl-Down moves down a line as Down does, and so on.
Appendix F lists every such key.

\WPSide{Function Keys}
Most InitechWord features are started with the function keys.  This book
names a feature in bold with its key in parentheses, as the template does:
\FK{Bold}{F6}, \FK{Format}{Shift-F8}.  A hyphen between two keys means hold
down the first while you press the second: \K{Shift-F8} means hold down
Shift and press F8.  A comma means release the first key before you press
the next: \KC{Shift-F8}{1} means press Shift-F8, then type 1.

The function keys InitechWord uses are these.  Appendix F lists them again
with every key of the keyboard.

\par\addvspace{4pt}
\noindent\begin{tabular}{@{}llp{52mm}@{}}
\textbf{Feature} & \textbf{Key} & \textbf{What it does}\\
Cancel & F1 & backs out of a menu or prompt; restores deleted text\\
Setup & Shift-F1 & units of measure and backup settings\\
\rarr Search & F2 & finds text and codes, forward\\
\larr Search & Shift-F2 & finds text and codes, backward\\
Replace & Alt-F2 & replaces text and codes\\
Reveal Codes & Alt-F3 or F11 & shows the hidden codes\\
\rarr Indent & F4 & indents a paragraph to the next tab stop\\
Block & Alt-F4 or F12 & marks text to act on\\
Move & Ctrl-F4 & moves, copies or deletes text\\
Bold & F6 & bold text\\
Exit & F7 & leaves the document\\
Underline & F8 & underlined text\\
Format & Shift-F8 & margins, justification, line spacing\\
Font & Ctrl-F8 & size and appearance of text\\
Save & F10 & saves the document\\
Retrieve & Shift-F10 & brings a saved document back\\
\end{tabular}
\par\addvspace{4pt}

Most other function keys name a WordPerfect feature that InitechWord does
not carry out: Shift-F7, for example, is Print.  Pressing one shows the
feature's name in grey on the status line, as on the template, and changes
nothing; the next key you press clears the name and does what it always
does.  Shift, Alt and Ctrl with F11 and F12 do nothing at all.  Appendix F
lists every function key.

\WPSide{Pull-Down Menus}
InitechWord has one pull-down menu, \textbf{File}, in the second row of the
menu bar.  It holds the three things you do with a whole document:
\textbf{Retrieve}, \textbf{Save} and \textbf{Exit}.  Point to \textbf{File}
with the mouse, hold down the mouse button, drag to the item, and release
the button.  Each item starts the same feature as its function key:
Retrieve as \K{Shift-F10}, Save as \K{F10}, Exit as \K{F7}.  The menu shows
the keyboard equivalent of Exit, the command key and Q: on your keyboard,
hold down Ctrl and press Q.

In this book, a step that can also be carried out from the File menu is
followed by the menu route in italic, after a small picture of the menu bar:
\PullDown{Select \Mi{S}ave from the \Mi{F}ile menu.}
For everything else you use the keys.  See \InRef{Pull-Down Menus}.

\WPSide{Menu Options}
When you press a feature key, InitechWord shows its menu on the status line
or, for Format and Setup, on a screen of its own.  Each option has a number.  Type
the number to choose the option.  This book names the option with its
number in parentheses: select \Ch{Size}{1}.  The \Scr{0} at the end of a
status-line menu is what you get if you choose nothing: press \FK{Cancel}{F1} or
any key that is not an option, and you are back in the document.

A menu shows the whole set of options WordPerfect users know, and
InitechWord carries out the options described in this book.  Where a
menu shows an option this book does not describe, pressing its number
returns you to the document (on the status-line menus) or leaves the
screen as it is (on the Format and Setup screens, where such options are
printed in grey on Setup).  The Font menu, for example,
shows Base Font (4) and Print Color (5); either returns you to the
document.

When InitechWord asks a question, the answer you get by pressing
\K{Enter} or any key other than the one shown in parentheses is the one
\emph{not} in parentheses.  At \Scr{Exit WP? No (Yes)}, type \Typed{y} to
leave InitechWord; any other key answers No.

\WPSide{Codes}
InitechWord keeps the formatting of a document as hidden codes in the text:
a code where bold begins and another where it ends, a code for each hard
return, tab, margin setting and so on.  The editing screen shows you their
effect -- bold text, an indented paragraph -- but not the codes themselves.
To see them, press \FK{Reveal Codes}{Alt-F3}: the screen divides, and the
lower part shows the text with every code in square brackets, \Code{[BOLD]}
where bold begins and \Code{[bold]} where it ends, \Code{[HRt]} for a hard
return.  Press Reveal Codes again to return to the full editing screen.

If you erase a code by accident on the editing screen, InitechWord first
asks you whether you mean to.  See \InRef{Codes} and \InRef{Reveal Codes};
Appendix B lists every code.

\WPSide{Exit}
When you have finished with a document, press \FK{Exit}{F7}.  InitechWord
always asks first whether you want to save the document, then whether you
want to leave InitechWord: answering No to the second question clears the
screen for a new document.  See \InRef{Exit}.
